% BEGIN PROFESSOR FEEDBACK PF13: Method - Evaluation
\section{Evaluation}
\label{sec:evaluation_protocol}

\subsection{Scope}

The evaluation measures how well an address can be classified after the fact, from
its full recorded history. It does not estimate how early an address could have
been caught, and it makes no causal claim about how a fraud was carried out. As
described in Section~\ref{sec:system_overview}, the data and every setting are
fixed in a preparation step before any model is trained, so nothing in the
evaluation can be tuned on the test data. The parts that enforce this, such as the
recorded settings and the checks that the code and data have not changed, are
described in Section~\ref{sec:implementation}.

\subsection{Cohorts and splits}

RQ1 and RQ2 are evaluated on three non-overlapping samples from the HuggingFace
dataset. Each sample contains two separate groups of addresses. An address is
eligible only if it takes part in at least three transactions, since
fewer cannot form one of the three-edge motifs. Before the groups are drawn, the
most extreme addresses are removed: within each combination of label, contract
type, and activity level, the top 0.5\,\% most active addresses are dropped.
Without this, a few exchange-scale hubs would dominate the graph on their own. The
cut is made within each combination rather than across the whole dataset, because
the busiest addresses are almost all benign exchange wallets, so a single global
cut would remove far more benign than fraudulent addresses and quietly shift the
balance between the classes. The rule is fixed before any model is trained, and
every removed address is recorded.

The two groups serve different purposes. The \emph{prevalence} cohort is a plain
random sample of up to 5{,}000 eligible addresses, which keeps the natural mix of
fraudulent and benign addresses and checks that the findings hold on a realistic
distribution. The \emph{controlled} cohort is drawn from the remaining addresses,
up to 250 per combination of label, contract type, and activity level, which
balances the classes and holds activity and contract type roughly constant so that
a difference cannot be explained by those alone. The two cohorts within a
sample share no addresses and are reported separately. After both cohorts are
drawn, their addresses and known campaign groups are removed from the eligible
pool before the next sample is created. Repeating this procedure produces three
samples without reusing addresses or known campaigns.

Splitting the addresses needs one precaution. Some fraud arrives in campaigns of
many near-identical addresses, and if some of them landed in the training set
while their near-twins landed in the test set, the detector would look better than
it really is. Campaign groups are therefore kept together, and every split keeps
a whole group on a single side. In the implementation, these groups are defined
from the published address descriptions. After trimming whitespace and converting
the text to lower case, descriptions beginning with
\emph{multichain hack alert}, \emph{upbit hacker}, or
\emph{sorbet finance hack alert} are assigned to their respective shared
family. Every other address receives its own group. These are description-based
grouping rules, not a complete reconstruction of coordinated campaigns.
Each cohort is divided about
80 to 20 into a development set and a test set with a group-aware, stratified
split, and the development set is divided again into repeated five-fold splits,
under different random seeds, for choosing the time window. Cohort membership is
included in the stratification, while campaign assignments are made across both
cohorts together. A campaign represented in both cohorts therefore remains on
the same side and in the same validation fold. RQ3 and RQ4 share their own 80 to
20 split, stratified by category, again keeping only categories with at least 100
addresses and never merging them.

\subsection{Metrics and window selection}

The main measure is average precision (AP), which summarises how well the detector
ranks malicious addresses above benign ones across all decision thresholds. It
suits imbalanced data because it focuses on the rarer positive class. The F1
score is the harmonic mean of precision and recall. These three measures are
reported as secondary measures at a fixed threshold of
0.5, which is never tuned on the test set. The area under the receiver operating
characteristic curve (ROC-AUC) is reported alongside them.
For RQ3 and RQ4 the matching measures are the macro-averaged and per-category
average precision, the macro F1, and a confusion matrix.

The time window is chosen on the development data alone. Each window is scored by
cross-validation, and the window with the highest average score is taken. If
several windows share that score, the shortest of them is used, so that the choice
is decided in advance rather than by chance. As a check on how firm that choice is,
the shortest window whose score is still within one standard error of the best is
recorded alongside it as the one-standard-error (1-SE) choice, which shows whether a shorter window would have done
essentially as well. The test set plays no part in either.

\subsection{Statistical tests}

For RQ1 the effect is the difference in test-set average precision between the
combined view and the baseline. Two procedures support it. The first is a
paired permutation test. For each campaign group, a random choice determines
whether the baseline and combined prediction scores are exchanged for every
address in that group. Address labels, campaign membership, and the train and
test assignments stay fixed. The models are not retrained. The difference in
AP is recomputed for 10{,}000 such exchanges, giving a one-sided $p$-value under
the null that the two sets of prediction scores are exchangeable within groups.

The second procedure is a bootstrap with 10{,}000 repetitions. Groups are
sampled with replacement separately within the benign and malicious classes.
Each class retains its original number of sampled groups. Every sampled group
contributes all of its addresses, including repeated copies if the group is
drawn more than once. The paired predictions are kept together and the AP
difference is recomputed, giving a 95\,\% confidence interval (CI).
The groups are the description-based families and singleton addresses defined
above.

Groups can contain different numbers of addresses. Each group has the same
chance of being sampled within its class, but a larger sampled group contributes
more addresses to AP. The total number of addresses and the class proportions
can therefore vary between bootstrap repetitions. The score-swap test, in
contrast, always uses the original test addresses. Both procedures preserve
dependence within the known groups, while the reported metric remains an
address-level measure. In the binary Real-CATS check, groups are individual
addresses because campaign identifiers are unavailable.

For RQ2 the chosen window is compared with each of the other windows using the
same kind of bootstrap. Because many windows are compared at once, the resulting
$p$-values are adjusted for multiple testing with the Benjamini--Hochberg
procedure (BH). The development curve and fingerprint coverage describe how
usefulness changes with $\delta$. The pairwise tests additionally check whether
the selected window is distinguishable from the alternatives. A unique best
window within the tested grid is supported only if it beats every alternative
after correction.

RQ3 asks whether the fingerprints contain category information. The held-out
category labels are shuffled 10{,}000 times while the fitted model and its
prediction probabilities stay fixed. Each shuffle preserves the class sizes.
The resulting macro AP values form a reference for no association between the
predictions and the category labels. Failure to reject this null would mean
insufficient evidence of category separation in this test. It would not prove
that the fingerprints contain no category information.

RQ4 asks whether the categories are better served by their own windows than by one
shared window. The two settings are compared on the test set in two ways. The first
is an overall comparison: the model using each category's own window is set against
the model using the single global window, and the difference in macro-averaged
average precision is measured with a paired bootstrap. The second looks at each
category on its own, comparing its own window with the global one, and because this
is repeated for every category the results are adjusted for multiple testing. The
number of distinct windows the categories select is reported alongside, since
category-specific windows can only help if the categories do not all choose the
same one.

\subsection{Robustness to event ordering}

Every result is produced under each of the seven orderings of
Section~\ref{sec:motif_extraction}. An ordering-dependent conclusion, such as an
RQ1 improvement or an RQ2 best window, is accepted only when its direction and its
decision agree across all seven orderings. This shows that a conclusion does not
depend on how transfers sharing a timestamp are ordered. It does not recover the
real order of those transfers, which the data does not record in the first place.
% END PROFESSOR FEEDBACK PF13
